\chapter{Острови, хаос і помилка \texorpdfstring{$\sqrt q$}{√q}}\label{ch:g08}

Це найповчальніша історія проєкту. Хибний канонічний імпульс ($\psi$ замість $\psit$) дав
формулу ширини магнітного острова, завищену в $\sqrt q$ раз. Вимірювання чесно показали
«дефіцит» ширини, а ми двічі пояснили його фізикою: спершу дискретністю засіву, потім
стохастичним шаром сепаратриси. Лише коли гамільтоніан силової лінії вивели заново, дефіцит
зник: відношення «виміряно/теорія» стали $0{,}95$--$1{,}02$. Помилка не нова: це рецидив R4
батьківського реєстру \cite{RegEst}, який журнал верифікації впіймав ще 22.09 у чернетці
тексту (V-D3 \cite{RegVerif}). Тоді її виправили в документі, але не в коді. Реєстри
\cite{RegViz} зберігають усі проміжні пояснення як спростовані (VR12, VR15, VR16). Прочитайте
цей розділ як детектив: кожен крок у ньому був розумним.

\section{Стенд: навмисні острови в реальній рівновазі}\label{sec:g08-stand}

Усе рахується на одному кадрі DIII-D 203702 (кадр 75, $t=\SI{1760}{мс}$, $\qnf=3{,}6995$). В
осесиметрії силові лінії інтегровні, і перетин Пуанкаре дає рівно контури $\psi$ (E18). Тому
острови треба \emph{внести} збуренням потоку \cite{g08Physics}:
\begin{equation}\label{eq:g08-pert}
  \psi_{\text{total}}=\psi(R,Z)+\sum_k A_k\,g_k(\psiN)\cos\!\left(m_k\thetastar-n_k\varphi+\alpha_k\right),
  \qquad \Bvec=\grad\times(\psi_{\text{total}}\grad\varphi)+F\grad\varphi,
\end{equation}
де $g=\psiN^{m/2}$~--- «вакуумна» обвідна (розд.~\ref{ch:g09} покаже, що для реальних котушок
вона не годиться, VR14). Поле є ротором, тож $\grad\cdot\Bvec=0$ точно для будь-якого числа мод
(VE21, розд.~\ref{ch:g07}).

Перш ніж шукати острови, стенд пройшов перевірку ланцюга $\psi\to\Bvec$ (VE1--VE7). Пошук осі
збігається з EFIT до $\SI{0,0002}{мм}$ (VE1). LCFS набору є контуром $\psi$ з розкидом 0,022\,\%
(VE2). Калібрування $F$ лише за $\qnf$ дає $\Btor(\text{вісь})=\SI{1,927}{Тл}$ за реальних
1,9--2,1~Тл (VE3). Два способи обчислити $q$ збігаються до 0,001--0,16\,\% (VE4). Інтегратор
DOP853 зберігає $\psi$ до $10^{-7}$ за 200 обертів (VE5). Аналітичні похідні збурення збігаються
зі скінченними різницями до $10^{-7}$ (VE7). Дві ранні помилки варто запам'ятати. Резонанс
вимагає кута прямих силових ліній $\thetastar$, а не геометричного (VR1). Растеризація
$\cos m\thetastar$ на сітку $65\times65$ давала 2,4\,\% паразитної $m=1$ і фальшивий ланцюг на
$q=1$ (VE6, VR2). Ширина острова йде як корінь з амплітуди, тож кілька відсотків «бруду» дають
добре видимий острів. З двома модами 2/1 і 3/1 при амплітуді $10^{-3}$ розріз показує
\textbf{2 острови на $\psiN=0{,}752$ ($q=2$) і 3 на $\psiN=0{,}900$ ($q=3$)} (VE8).

\section{Як розпізнати острів: ланцюг спростувань класифікатора}\label{sec:g08-classifier}

Щоб порівнювати ширину з теорією, орбіти треба автоматично розкласти на регулярні, острівні
й хаотичні. Кожен шар класифікатора в \texttt{tokviz/analysis.py} з'явився після спростування
попереднього (VR8--VR13).

\paragraph{Крок 1. Число обертання (VR8 $\to$ VE16).} Просте середнє
$\nu=\langle\dd\pol/\dd\varphi\rangle$ за скінченним вікном несе осциляційну похибку
$\order{1/N}$, бо $\pol$ наростає вздовж D-подібної поверхні нерівномірно. На
\emph{інтегровному} полі це дало 50\,\% «хаотичних» орбіт (VR8). Заміна~--- зважене усереднення
Біркгофа з гладкою «шапочкою» $w(t)=\exp[-1/(t(1-t))]$. Вага разом з усіма похідними
занулюється на кінцях, тож на квазіперіодичній орбіті середнє збігається надполіноміально, а на
хаотичній не збігається зовсім (у докстрингу коду названо Das \& Yorke і Sander \& Meiss).
Так виміряне $\nu$ відтворює $1/q$ рівноваги до 0,03\,\%: це третій незалежний спосіб
дістати $q$ (VE16).

\code{viz/src/tokviz/analysis.py}{42}{57}{(ваги й зважене середнє Біркгофа)}

Рядки 42--51: вузли $t_i=(i+\frac12)/n$ лежать усередині $(0,1)$, тож $w$ ніде не ділить на
нуль. Нормування $\sum w_i=1$ робить результат справжнім середнім. Рядки 54--57: середнє вздовж
осі~1 (вздовж орбіти) рахується одним матричним множенням для всіх орбіт одразу.

\code{viz/src/tokviz/analysis.py}{81}{96}{(збіжність у десяткових цифрах)}

Рядки 88--92: $\nu$ за всією орбітою і за її першою половиною. Рядки 93--96: кількість
цифр, у яких вони збігаються, $-\log_{10}|\nu_{\text{full}}-\nu_{\text{half}}|/\nu_{\text{full}}$.
Поріг хаосу не вгадують, а калібрують на $A=0$, де система точно інтегровна: мінімум
5,14 цифри, медіана 6,66, поріг на декаду нижче, $4{,}14$ \cite{g08Analysis}.

\paragraph{Крок 2. Острів --- це захоплення $\nu$, а не екскурсія (VR9--VR13).} Перша ідея
«острівна орбіта має велику радіальну екскурсію $\psiN$» хибна в принципі (VR9). У фігурній
рівновазі екскурсія зростає до краю й для регулярних орбіт, а справжні острівні орбіти
($\nu=0{,}49999$) класифікувались як хаос. Правильна ознака~--- \emph{фазове захоплення}: у
смузі острова всі орбіти мають однакове $\nu=n/m$. Далі йде ланцюжок уточнень, і кожне
спростоване нульовим збуренням, де островів бути не може:
\begin{itemize}
  \item близькості до раціонального числа у вікні 2\,\% замало: на $A=0$ це дало 2 «острови» (VR10);
  \item додали взаємну узгодженість: $\nu$ членів смуги збігаються між собою до $2\cdot10^{-3}$;
  \item при 150 зернах сусіди мають майже однакове $\nu$ просто через гладкість профілю,
        і знову вийшло 2 «острови» на $A=0$ (VR13);
  \item додали безрозмірний \emph{тест пласкості}: розкид $\nu$ у смузі порівнюють зі
        зміною рівноважного $1/q$ на тому самому інтервалі $\psiN$.
\end{itemize}

\code{viz/src/tokviz/analysis.py}{133}{160}{(\texttt{find\_locked\_bands}: перевірка однієї смуги)}

Рядки 136--140: середнє $\nu$ смуги та її відносний розкид (тест VR10). Рядки 142--152: тест
пласкості (VR13). \texttt{expected}~--- на скільки змінилось би $\nu$ на ширині смуги без
захоплення. Тест вирішальний лише тоді, коли ця зміна помітна (рядок~151). Рядки 154--160:
смуга з кількох орбіт або одна дуже точно захоплена орбіта (\texttt{solo\_tol}).

\paragraph{Крок 3. Порядок перевірок (VR11).} Острівна орбіта має другу, повільну частоту~---
лібрацію навколо O-точки~--- і тому за Біркгофом збігається \emph{гірше} за KAM-орбіту. Якщо
спершу перевірити збіжність, острів завжди читається як хаос. Тому захоплення перевіряється
першим:

\code{viz/src/tokviz/analysis.py}{193}{221}{(тіло \texttt{classify})}

Рядки 193--200: цифри збіжності, $\nu$ і смуги захоплення. Рядки 202--208: острівними стають
члени смуг, які вижили й мають хоч якусь збіжність (\texttt{stability\_digits}). Рядок 210:
хаос~--- лише серед \emph{не}острівних орбіт. Рядки 213--221: мітки й частка хаосу.
Результат видно на рис.~\ref{fig:g08-ab}: хаотичні орбіти сидять рівно на $\nu=1/2$ і $\nu=1/3$,
хоча класифікатор нічого не знає про резонанси (VE17).

\begin{figure}[ht]
  \centering
  \includegraphics[width=\textwidth]{g08_poincare_after_AB.jpg}
  \caption{Класифікований перетин при $A=10^{-3}$ (A) і профіль $\nu(\psiN)$ проти $1/q$
  рівноваги (B), DIII-D 203702. Свіжий прогін 2026-09-29 з виправленою формулою. Дві
  помаранчеві орбіти на $\nu=2/3$ ($\psiN\approx0{,}61$)~--- смуга 3/2, яку класифікатор
  знаходить і при $A=0$ (за \texttt{data/analysis/analysis.json}). Це прояв VQ4
  (розд.~\ref{sec:g08-vq4}). Власний розрахунок, код: \texttt{tokamak-3d-viz/src/run\_analysis.py}.}
  \label{fig:g08-ab}
\end{figure}

\section{Ширина острова: маятник і канонічна пара}\label{sec:g08-pendulum}

\begin{derivation}[ширина острова з гамільтоніана силової лінії]
Для $\Bvec=\grad\psi\times\grad\varphi+F\grad\varphi$ ($\psi$~--- полоїдальний потік на радіан)
рівняння силових ліній є гамільтоновою системою з 1,5 ступеня вільності (E19,
\cite{Escande2024}): час~--- $\varphi$, координата~--- $\thetastar$,
\textbf{імпульс~--- тороїдальний потік $\psit$}, \textbf{гамільтоніан~--- полоїдальний потік $\psi$}:
\begin{equation}\label{eq:g08-ham}
\frac{\dd\thetastar}{\dd\varphi}=\frac{\pa\psi}{\pa\psit}=\frac1q\equiv\nu,\qquad
\frac{\dd\psit}{\dd\varphi}=-\frac{\pa\psi}{\pa\thetastar},\qquad \dd\psi=\frac{\dd\psit}{q}.
\end{equation}
Нехай $\psi=\psi_0(\psit)+A\cos\xi$ з резонансною фазою $\xi=m\thetastar-n\varphi+\alpha$. Біля
резонансу $\nu(\psit^{\text{res}})=n/m$ розкладаємо $\nu$ до першого порядку й дістаємо маятник:
\begin{equation}\label{eq:g08-pend}
\frac{\dd\xi}{\dd\varphi}=m\,\nu'\,(\psit-\psit^{\text{res}}),\qquad
\frac{\dd\psit}{\dd\varphi}=mA\sin\xi,\qquad
\nu'=\frac{\dd\nu}{\dd\psit}=-\frac{1}{q^2}\frac{\dd q}{\dd\psi}\frac{\dd\psi}{\dd\psit}=-\frac{\dd q/\dd\psi}{q^3}.
\end{equation}
Це рівняння руху для $K=\frac12 m\nu' p^2+mA\cos\xi$, $p=\psit-\psit^{\text{res}}$. На
сепаратрисі $\frac12|m\nu'|p_{\max}^2=2|mA|$, звідки повна ширина $4\sqrt{A/|\nu'|}$ (множник $m$
скорочується). Ширина в $\psi$ у $q$ разів менша за ширину в $\psit$:
\begin{equation}\label{eq:g08-W}
W_{\psi}=\frac1q\cdot4\sqrt{\frac{Aq^3}{|\dd q/\dd\psi|}}=4\sqrt{\frac{Aq}{|\dd q/\dd\psi|}}
\quad\Longrightarrow\quad
W=4\sqrt{\frac{\varepsilon q}{|\dd q/\dd\psiN|}},\qquad \varepsilon=\frac{A}{|\psib-\psiax|}.
\end{equation}
Друга форма~--- та сама формула в нормованому потоці: $A$ і $\dd q/\dd\psi$ масштабуються з
діапазоном потоку однаково.

\emph{Де була помилка.} Якщо вважати імпульсом сам $\psi$, тобто писати
$\dd\psi/\dd\varphi=mA\sin\xi$, то той самий маятник дає
$W=4\sqrt{A/|\dd\nu/\dd\psi|}=4\sqrt{Aq^2/|\dd q/\dd\psi|}$. Насправді
$\dd\psi/\dd\varphi=\nu\,\dd\psit/\dd\varphi$: у рівнянні для $\dot\psi$ бракує множника $1/q$, і
ширина виходить завищеною в $\sqrt q$ раз. Для $q=2$ це 41\,\%, для $q=3$~--- 73\,\%.
\end{derivation}

Сьогодні формула в коді така:

\code{viz/src/tokviz/perturbation.py}{197}{230}{(\texttt{island\_width\_psin}, виправлена версія)}

Рядки 198--224: докстринг повторює виведення~\eqref{eq:g08-ham}--\eqref{eq:g08-W}, щоб наступний
читач не мусив його відновлювати. Рядки 225--226: запис про стару формулу (правило «нічого не
зникає мовчки» діє й у коді). Рядки 228--230: власне обчислення; до 2026-09-29 тут стояло
\verb|eps*q**2| \cite{g08Changes}.

Перекриття двох ланцюгів оцінює параметр Чирикова
$S=(W_1+W_2)/(2|\psi_{N,1}-\psi_{N,2}|)$ (функція \texttt{chirikov}, рядки 233--242 того самого
файлу; у докстрингу прямо сказано, що $S\ge1$~--- евристика, а не теорема). Для стенда при
$A=10^{-3}$: мода 2/1 має $\psiN=0{,}7520$, $\varepsilon=7{,}52\cdot10^{-4}$,
$\dd q/\dd\psiN=4{,}44$, $W=0{,}0736$. Мода 3/1 має $\psiN=0{,}9001$,
$\varepsilon=8{,}54\cdot10^{-4}$, $\dd q/\dd\psiN=10{,}70$, $W=0{,}0619$. Звідси
$S=\mathbf{0{,}458}$ (VE9; до виправлення $W=0{,}1041/0{,}1072$ і $S=0{,}714$).

\section{Історія хибного пояснення}\label{sec:g08-story}

\paragraph{Вимір.} Ширину острова міряє \texttt{measured\_island\_width}: це об'єднання
радіальних екскурсій орбіт, захоплених на резонансі, тобто смуга, яку перетин справді показує.

\code{viz/src/tokviz/analysis.py}{224}{241}{(виміряна ширина ланцюга)}

Рядки 231--235: члени смуг саме цього резонансу з міткою «острів». Рядки 238--241: ширина~---
розмах $\psiN$ уздовж \emph{траєкторій} (а не лише початкових зерен), плюс центр і кількість
орбіт.

\begin{refuted}{VR12}
«Дефіцит виміряної ширини пояснюється дискретністю засіву.» При 56 зернах (крок 0,0215)
поправка $W_{\text{теор}}-2\Delta\psi$ чисельно збігалася з виміром, і це виглядало переконливо.
Після згущення до кроку 0,003--0,004 біля резонансів поправка стала нікчемною, а відношення
0,72 лишилося. Збіг при 56 зернах був випадковим.
\end{refuted}

\begin{refuted}{VR16}
«Дефіцит (0,59--0,74)~--- це стохастичний шар сепаратриси, що з'їдає регулярне ядро острова.»
Це була \emph{опублікована} інтерпретація у VE18 і «справжня причина» у VR12. Розріз упоперек
$q=2$ справді показує тонкі хаотичні шари обабіч захопленого ядра $\psiN\in[0{,}7429;\,0{,}7739]$
з $\nu=0{,}50000$. Але пояснення не витримує двох перевірок. Відношення не залежало від
амплітуди, хоча шар, що росте з $S$, мав би давати таку залежність. Зате для кожної моди воно
збігалося з $1/\sqrt q$: $1/\sqrt2=0{,}71$, $1/\sqrt3=0{,}58$.
\end{refuted}

Саме останнє спостереження мало насторожити: безрозмірне число, яке залежить лише від $q$
резонансу,~--- підпис помилки у формулі, а не фізики. Помилку знайшов координатор методички під
час написання розд.~2 методички (знахідка N1 \cite{g08Coord}). Формулу виправили 2026-09-29.

\begin{refuted}{VR15}
«Ширину острова можна рахувати, беручи $\psi$ за канонічний імпульс»:
$W=4\sqrt{\varepsilon q^2/|\dd q/\dd\psi|}$. Це повторення R4 батьківського реєстру.
Канонічна пара~--- $(\thetastar,\psit)$, гамільтоніан~--- $\psi$ (E19). Виправлено в
\texttt{perturbation.island\_width\_psin} і \texttt{coilfield.CoilPerturbation.calibrate\_to\_island\_width}.
\emph{Не} оновлено: \texttt{data/bake01} і рендери Blender досі несуть старі $W$ і $S$.
\end{refuted}

Ті самі вимірювання (вихідний засів, збережений у \texttt{data/analysis\_pre\_fix\_2026-09-29/})
переоцінено новою формулою. Самі вимірювання від формули не залежать
(табл.~\ref{tab:g08-ratios}, рис.~\ref{fig:g08-cd}).

\begin{table}[ht]
\centering\footnotesize\setlength{\tabcolsep}{4pt}
\caption{Виміряна ширина острова проти маятникової: до і після виправлення
(\texttt{docs/ANALYSIS.md} §2.2; звірено з обома файлами \texttt{analysis.json})}
\label{tab:g08-ratios}
\begin{tabular}{@{}llcccccc@{}}
\toprule
$A$ & мода & $W_{\text{вим}}$ & $W_{\text{стар}}$ & вим./стар. & $W_{\text{нов}}$ & вим./нов. & $S$: стар.\,$\to$\,нов.\\
\midrule
$5\cdot10^{-4}$ & 2/1 & 0,0492 & 0,0736 & 0,67 & 0,0520 & \textbf{0,95} & 0,505\,$\to$\,0,324\\
$10^{-3}$ & 2/1 & 0,0752 & 0,1041 & 0,72 & 0,0736 & \textbf{1,02} & 0,714\,$\to$\,0,458\\
$10^{-3}$ & 3/1 & 0,0628 & 0,1072 & 0,59 & 0,0619 & \textbf{1,01} & 0,714\,$\to$\,0,458\\
$2\cdot10^{-3}$ & 2/1 & 0,1044 & 0,1472 & 0,71 & 0,1041 & \textbf{1,00} & 1,009\,$\to$\,0,647\\
$4\cdot10^{-3}$ & 3/1 & 0,1578 & 0,2145 & 0,74 & 0,1238 & 1,27$^{*}$ & 1,427\,$\to$\,0,915\\
\midrule
\multicolumn{8}{@{}l}{Трасування сепаратриси без класифікатора, одна мода, $A=10^{-3}$:
\textbf{0,983} (2/1), \textbf{0,996} (3/1) від нової формули}\\
\bottomrule
\end{tabular}\\[2pt]
{\footnotesize $^{*}$Поза маятниковим режимом: острів 3/1 сягає $\psiN\approx0{,}96$, $S\approx0{,}9$,
смуга з двох орбіт.}
\end{table}

\begin{figure}[ht]
  \centering
  \begin{subfigure}{0.82\textwidth}\includegraphics[width=\linewidth]{g08_poincare_before_CD.jpg}
  \caption{до виправлення: $W=4\sqrt{\varepsilon q^2/|\dd q/\dd\psi|}$}\end{subfigure}\\[3pt]
  \begin{subfigure}{0.82\textwidth}\includegraphics[width=\linewidth]{g08_poincare_after_CD.jpg}
  \caption{після: $W=4\sqrt{\varepsilon q/|\dd q/\dd\psiN|}$ (свіжий прогін)}\end{subfigure}
  \caption{Ширина острова проти амплітуди (ліворуч; лінії~--- теорія, кружки~--- вимір) і частка
  хаотичних орбіт проти $S$ Чирикова (праворуч). До виправлення всі виміряні точки лежать нижче
  теорії на сталий множник $\approx1/\sqrt q$. Після виправлення вони лягають на лінії, а скан
  уже не досягає $S=1$. У свіжому прогоні класифікатор не знайшов смугу 2/1 при $5\cdot10^{-4}$ і
  $2\cdot10^{-3}$ (VQ4). Власний розрахунок, код: \texttt{tokamak-3d-viz/src/run\_analysis.py}
  (\texttt{data/analysis\_pre\_fix\_2026-09-29/} і \texttt{data/analysis/}).}
  \label{fig:g08-cd}
\end{figure}

Вирішальним став незалежний тест без класифікатора. Він трасує одну моду й вважає острівною
орбіту, у якої резонансна фаза $\xi=m\thetastar-n\varphi$ лібрує (розмах менший за $2\pi$):

\code{viz/tests/test_physics.py}{139}{161}{(\texttt{test\_island\_width\_matches\_traced\_separatrix})}

Рядки 140--145: одна мода, її резонанс, $\varepsilon$ і передбачена $W$. Рядки 147--150:
20 зерен на зовнішній середній площині в межах $\pm0{,}75W$ і трасування на 50 обертів. Рядки
151--153: критерій лібрації, незалежний від $\nu$ і від класифікатора. Рядок 155: зерна мусять
охоплювати острів з обох боків. Рядки 158--161: вимір має бути в межах 10\,\% від нової
формули і відрізнятися від старої більш ніж на 20\,\%. Отже, тест здатен розрізнити
$\sqrt q$. Повторний запуск цього коду для путівника (2026-09-29) дав $W_{\text{вим}}=0{,}0723$
проти $0{,}0736$ (2/1, відношення $0{,}983$, 8 захоплених орбіт) і $0{,}0617$ проти $0{,}0619$
(3/1, $0{,}996$, 5 орбіт); реєстрові числа відтворились.

\begin{established}{VE18 (переформульовано 2026-09-29)}
Маятникова формула (виправлена) справджується: виміряна ширина збігається з передбаченою до
0--5\,\% при $A\le2\cdot10^{-3}$. Тонкі стохастичні шари біля сепаратриси є, але траєкторії
захоплених орбіт (0,0752) заповнюють усю передбачену сепаратрису (0,0736). «Три ширини» старого
тексту зводяться до двох: розкид \emph{зерен} ядра (0,0310; зерна стоять на одній лінії при
одній фазі) і ширина сепаратриси.
\end{established}

\section{Хаос і критерій Чирикова}\label{sec:g08-chirikov}

Виправлення змінило вісь $S$, але не частки хаосу: їх класифікатор міряв незалежно від формули
(табл.~\ref{tab:g08-chaos}).

\begin{table}[ht]
\centering\small
\caption{Частка хаотичних орбіт (VE19, VE20, VE33; \texttt{docs/ANALYSIS.md} §2.3)}
\label{tab:g08-chaos}
\begin{tabular}{@{}lccccc@{}}
\toprule
$A$ & $S$ стар.\,$\to$\,нов. & незміщена (80 рівномірних) & усі 150 зерен & завищення & свіжий прогін\\
\midrule
0 & 0 & 0,000 & 0,000 & --- & 0,000\\
$2{,}5\cdot10^{-4}$ & 0,357\,$\to$\,0,229 & 0,075 & 0,133 & 1,78$\times$ & 0,075\\
$5\cdot10^{-4}$ & 0,505\,$\to$\,0,324 & 0,125 & 0,200 & 1,60$\times$ & 0,100\\
$10^{-3}$ & 0,714\,$\to$\,0,458 & 0,113 & 0,213 & 1,90$\times$ & 0,125\\
$2\cdot10^{-3}$ & 1,009\,$\to$\,0,647 & 0,175 & 0,300 & 1,71$\times$ & 0,200\\
$4\cdot10^{-3}$ & 1,427\,$\to$\,0,915 & 0,263 & 0,413 & 1,57$\times$ & 0,200\\
\bottomrule
\end{tabular}
\end{table}

\begin{itemize}
\item \textbf{VE9.} $S=0{,}458<1$, і поверхні між ланцюгами справді виживають. Це один
      узгоджений випадок, а не доказ евристики. Після виправлення висновок лише посилився.
\item \textbf{VE19.} Хаос існує задовго до $S=1$: вимірна частка є вже при $S=0{,}229$. Друга
      половина старого заголовка («різкого порогу при $S=1$ немає») знята: з виправленою формулою
      скан \emph{не досягає} $S=1$, тож перехід через одиницю цими даними не перевірено. Для
      $S\ge1$ потрібні $A\gtrsim5\cdot10^{-3}$.
\item \textbf{VE20.} Зерна, згущені біля резонансів, завищують частку хаосу в 1,6--1,9 раза.
      Тому статистику хаосу рахують лише на рівномірних 80 зернах, а згущені 70 йдуть тільки на
      вимір ширини. Методологічний урок: будь-яка «частка стохастичного об'єму», порахована на
      зернах біля резонансів, завищена приблизно вдвічі.
\item \textbf{VE33 (закрито VC5).} Немонотонність частки при 40 зернах (0,275\,$\to$\,0,200) була
      шумом вибірки. Одна орбіта з 80~--- це 1,25\,\%, тож западина 0,125\,$\to$\,0,113 незначуща.
      «Підтверджено» тут означає «у межах статистики однієї-двох орбіт».
\end{itemize}

Застереження, що стоїть на всіх числах: \emph{вакуумний} перетин переоцінює стохастичність, бо
плазма екранує резонансні компоненти. Частки в табл.~\ref{tab:g08-chaos}~--- верхня межа.

\section{Що лишилось відкритим: VQ4}\label{sec:g08-vq4}

Свіжий прогін тією самою командою змінив засів: згущення тепер ставиться за меншою
передбаченою шириною. Поріг хаосу став $3{,}88$ цифри замість $4{,}14$. Смугу 2/1 при $10^{-3}$
знайдено з тим самим відношенням $1{,}02$. Але при $5\cdot10^{-4}$ і $2\cdot10^{-3}$ класифікатор
\emph{не виділив} жодної захопленої смуги 2/1 (при $2{,}5\cdot10^{-4}$~--- лише плато), а при $A=0$
позначив 2 орбіти як «острови». Уточнення реєстру 2026-09-29: хибна \emph{смуга} одна~--- 3/2
(2 орбіти, розкид $\nu$ $1{,}2\cdot10^{-3}$), і вона з'являється при всіх $A\le10^{-3}$; 5/2 дає
лише $\nu$-плато шириною 0,0011, смугою не стає. Тобто тест пласкості, введений після VR13, сам
виявився чутливим до щільності засіву.

\begin{conjecture}{VQ4 (відкрите питання реєстру)}
Чи надійний класифікатор островів (тест пласкості $\nu$) при новому засіві? Висновки про ширину
(VE18) спираються на вихідний засів і незалежний тест лібрації, тож їх це не зачіпає. Проте
свіжий \texttt{analysis.json}~--- гірше джерело для VE19 і VE33, а \texttt{run\_parasitic}
після виправлення не перезапускався. \emph{Закриє:} тест пласкості, що на $A=0$ дає 0 островів
і знаходить 2/1 на всіх амплітудах при обох засівах.
\end{conjecture}

\begin{practicum}[відтворити й зламати]
Потрібні дані DIII-D (\texttt{TOKVIZ\_DATA\_DIR} на теку з \texttt{d3d\_shot\_203702.parquet}) і
Python зі scipy (підходить venv конкурсу \texttt{fusion equilibrium challenge/starter/.venv}).
\begin{enumerate}
\item \verb|python tests/test_physics.py| --- усі фізичні тести стенда, зокрема тест лібрації.
\item Поверніть у \texttt{island\_width\_psin} множник \verb|q_res**2| і запустіть тест знову. Він
      має впасти на рядку «measured \ldots vs predicted». Це і є регресійний захист від R4.
\item \verb|python src/run_analysis.py --out data/analysis --n-seeds 80 --n-refine 35| відтворює
      рис.~\ref{fig:g08-cd}. Додайте \verb|--amps 0,5e-3,1e-3,5e-3,8e-3|, щоб уперше перевірити
      поведінку при $S\ge1$ (відкрито після VE19).
\item Задача для VQ4: змініть \texttt{flatness} чи умову рядка~151 у \texttt{find\_locked\_bands}
      так, щоб при $A=0$ було 0 островів для \emph{обох} засівів. Засів задають
      \verb|--n-refine| і \verb|--refine-span|.
\end{enumerate}
\end{practicum}

\paragraph{Мораль.} Помилку в канонічному імпульсі проєкт уже раз спростував (R4, V-D3), і все ж
вона повернулася в коді. Її двічі «пояснили» фізикою, бо кожне пояснення правдоподібне локально.
Тривогу мали викликати два сигнали: безрозмірне відношення, яке не залежить від амплітуди, і його
збіг із простим числом ($1/\sqrt q$). Врятували її три речі: незалежний тест, реєстр, де старі
пояснення не стираються, і повторне виведення з перших принципів.
